\documentclass[11pt]{article}

\usepackage[margin=1.6cm]{geometry}
\usepackage{amsmath}
\usepackage{amssymb}
\usepackage{enumitem}
\usepackage{titlesec}
\usepackage{xcolor}
\usepackage{fancyhdr}
\usepackage{multicol}
\usepackage{array}

\titleformat{\section}{\large\bfseries\color{blue!50!black}}{\thesection}{0.6em}{}
\titlespacing{\section}{0pt}{1.1em}{0.4em}
\titleformat{\subsection}{\normalsize\bfseries}{\thesubsection}{0.6em}{}
\titlespacing{\subsection}{0pt}{0.8em}{0.3em}

\setlength{\headheight}{14pt}

\pagestyle{fancy}
\fancyhf{}
\fancyhead[L]{Statistical Questions and Variables}
\fancyfoot[C]{\thepage}
\renewcommand{\headrulewidth}{0.4pt}

\fancypagestyle{firstpage}{%
  \fancyhf{}
  \fancyhead[L]{Statistical Questions and Variables}
  \fancyhead[R]{Name: \makebox[4cm]{\hrulefill}}
  \fancyfoot[C]{\thepage}
  \renewcommand{\headrulewidth}{0.4pt}
}

\newcommand{\ans}{\hrulefill}
\newcommand{\anshere}{\par\vspace{0.4em}\noindent\ans\par\vspace{0.4em}}
\setlist[enumerate,1]{leftmargin=1.4em, itemsep=0.5em, topsep=0.3em}
\setlength{\columnsep}{1.1cm}
\setlength{\columnseprule}{0.4pt}
\raggedcolumns

\begin{document}
\thispagestyle{firstpage}

\begin{center}
  {\Large\bfseries Statistical Questions and Variables}
\end{center}
\vspace{0.3em}

\section{What is a Statistical Question?}

\noindent A \textbf{statistical question} is one that is answered using statistics --- by
\emph{collecting data} from a group (for example, asking lots of people, or measuring lots of
things). You expect the data you collect to \textbf{vary}, rather than there being one single
fixed answer.

\vspace{0.3em}
\noindent\textbf{Worked example:} Is ``How many pets does Lily have?'' a statistical question?
\emph{No} --- it asks about one specific person, so there is only one fixed answer. But ``How
many pets do students in Year 9 have?'' \emph{is} a statistical question, because if you asked
every student, you would expect the answers to vary from student to student, and you would need
to actually collect that data to find out.

\vspace{0.5em}
\noindent\textbf{For each question below, write S (statistical) or N (not statistical), and give
a one-sentence reason.}

\begin{multicols}{2}
\begin{enumerate}[label=(\alph*)]
  \item How many hours of screen time does Alex have per day?
  \anshere
  \item How many hours of screen time do students in our class have per day?
  \anshere
  \item What is the tallest building in the world?
  \anshere
  \item What are the heights of the buildings in Auckland's CBD?
  \anshere
  \item What is my dog's favourite treat?
  \anshere
  \item What are the favourite treats of dogs at the local vet clinic?
  \anshere
  \item How many points did the All Blacks score in their last match?
  \anshere
  \item How many points do rugby teams score, on average, in matches this season?
  \anshere
\end{enumerate}
\end{multicols}

\vspace{0.3em}
\noindent\textbf{Rewrite each question below so that it becomes a sound statistical question.}
\begin{enumerate}[label=(\alph*)]
  \item ``What's my favourite music genre?''
  \anshere
  \item ``How many hours did I sleep last night?''
  \anshere
  \item ``How fast can I run $100$\,m?''
  \anshere
\end{enumerate}

\newpage

\section{Explanatory vs Response Variables}

\noindent When a statistical question investigates a \textbf{relationship between two
variables}, we give those variables different names:
\begin{itemize}[leftmargin=1.4em, itemsep=0.2em, topsep=0.2em]
  \item The \textbf{explanatory variable} is the one we think might explain or influence a
    change --- it's the variable we control or choose.
  \item The \textbf{response variable} is the one being measured --- it may \emph{respond} to
    changes in the explanatory variable.
\end{itemize}
\noindent (This does not always mean one \emph{causes} the other --- only that we are
investigating whether they are related.)

\vspace{0.3em}
\noindent\textbf{Worked example:} Does the temperature outside affect how many ice creams a shop
sells?
\begin{itemize}[leftmargin=1.4em, itemsep=0.2em, topsep=0.2em]
  \item Explanatory variable: temperature.
  \item Response variable: number of ice creams sold.
\end{itemize}

\vspace{0.4em}
\noindent\textbf{For each scenario, state the explanatory variable and the response variable.}
\begin{enumerate}[label=(\alph*)]
  \item Does the amount of time spent practising free throws affect a basketball player's
    success rate?
  \anshere
  \item Is there a relationship between the number of hours of sleep someone gets and their
    test score the next day?
  \anshere
  \item Does the amount of rainfall in a month affect the height of plants in a garden?
  \anshere
  \item Researchers want to know if a person's daily step count affects their resting heart
    rate.
  \anshere
  \item Does the size of a classroom affect how many students can comfortably fit inside it?
  \anshere
  \item Is a person's typing accuracy related to how many years they have been using a
    keyboard?
  \anshere
\end{enumerate}

\newpage

\subsection*{Spot the mistake}
\noindent\textbf{For each statement, explain what is wrong with it.}
\begin{enumerate}[label=(\alph*)]
  \item A student investigating rainfall and plant height says: ``The explanatory variable is
    plant height and the response variable is rainfall.''
  \anshere
  \item A student investigating sleep and test scores says: ``The response variable is hours of
    sleep, because that's the one I measured first.''
  \anshere
\end{enumerate}

\vspace{0.5em}
\subsection*{Putting it all together}
\noindent\textbf{For each scenario: (i) decide whether it is a sound statistical question, and
(ii) if it investigates a relationship between two variables, state the explanatory and response
variables.}

\begin{enumerate}[label=(\alph*)]
  \item Do students who do more exercise each week tend to get better grades?
  \anshere
  \item What is the most common pet owned by households in New Zealand?
  \anshere
  \item Does wearing a seatbelt reduce the risk of injury in a car accident?
  \anshere
\end{enumerate}

\vspace{0.4em}
\noindent\textit{Hint for (b): some statistical questions only investigate \emph{one} variable
(like ``favourite pet''), rather than a relationship between two things --- these don't have an
explanatory or response variable at all.}

\vspace{0.6em}
\subsection*{Your turn}
\noindent Write your own statistical question that investigates a relationship between two
variables. Identify the explanatory variable and the response variable in your question.
\anshere
\anshere

\end{document}
